\documentclass[UTF8,11pt]{ctexart}
\usepackage[a4paper,margin=2.05cm]{geometry}
\usepackage{amsmath,amssymb,mathtools}
\usepackage{booktabs,longtable,tabularx,array}
\usepackage{xcolor}
\usepackage{enumitem}
\usepackage{hyperref}
\usepackage{fancyhdr}
\usepackage{lastpage}
\hypersetup{colorlinks=true,linkcolor=blue!55!black,urlcolor=blue!65!black,citecolor=blue!55!black}
\setlength{\parindent}{2em}
\setlength{\parskip}{0.35em}
\setlist{nosep,leftmargin=2em}
\pagestyle{fancy}
\fancyhf{}
\lhead{Swarm-EOD / MATCH 数学模型}
\rhead{原始 Gitee 提交 53e129b}
\cfoot{\thepage/\pageref{LastPage}}
\newcommand{\src}[1]{\textsuperscript{\textcolor{blue!65!black}{[#1]}}}
\newcommand{\clip}{\operatorname{clip}}
\newcommand{\wrap}{\operatorname{wrap}}
\newcommand{\E}{\mathbb{E}}
\newcommand{\Prob}{\mathbb{P}}
\newcommand{\R}{\mathbb{R}}
\newcommand{\one}{\mathbf{1}}
\newcommand{\given}{\,\middle|\,}

\title{MATCH / Swarm-EOD 红队任务完整数学模型\\
\large 基于原始 Gitee 规则与框架的独立重建}
\author{来源审计与数学形式化}
\date{2026年10月14日}

\begin{document}
\maketitle

\begin{center}
\fbox{\parbox{0.94\textwidth}{
\textbf{来源边界。}本文只依据 Gitee 仓库
\href{https://gitee.com/CJYLZS/swarm-eod}{CJYLZS/swarm-eod} 的固定提交
\texttt{53e129bd7d2efc82c29deac2ecd8c57a536e74e7}，读取规则文档、配置、模拟器和评分实现。
没有读取任何现有红队候选算法来推导状态、目标、约束或策略结构；当前候选的历史得分也没有进入模型。
仓库中的示例策略只视为示例，不视为规则。文中“源事实”“数学形式化”“近似模型”“待标定量”严格区分。
}}
\end{center}

\section{结论先行}

这个任务的精确类别是一个\textbf{有限时域、连续状态、混合动作、分散部分可观测的随机控制问题}，红队阶段结束后还要经过一个蓝队救援过程才得到完整终端回报。用控制论语言，它是连续状态的 Dec-POMDP；把蓝队视作给定参考策略时，是带二阶段随机终端奖励的 Dec-POMDP。只有在蓝队会针对红队布局重新优化时，才适合写成双层优化；原始评分实现并不要求这个更强假设。

默认红队总分的精确结构为\src{S1,S2}
\begin{equation}
S_R=0.5S_{\mathrm{map}}+0.2S_{\mathrm{bomb}}+0.2S_{\mathrm{wall}}+0.1S_{\mathrm{time}},
\qquad S_\bullet\in[0,100].
\label{eq:red-total}
\end{equation}
这里“覆盖范围分”是炸弹圆形作用区覆盖的\textbf{加权墙体}，即 $S_{\mathrm{wall}}$；无人机走过的面积只作为诊断量，默认权重为零，不能与地图精度或炸墙覆盖混为一谈。

原始资料没有规定“每张地图必须设置某几个随机种子”。种子是估计随机策略期望性能、做配对比较和复现实验的\textbf{评测设计}，不是游戏规则，也不是启发式控制目标。并且当前源码中 Playground 使用独立的 \texttt{default\_rng(seed)}，而多种传感器噪声直接使用全局 \texttt{np.random}；单独传入 Playground 种子并不能保证整个回合逐位复现。\src{S8,S9,S10}

分数长期没有明显改善的根本原因，可以从目标函数看清：四个分项由不同机制产生，且炸弹位置同时影响 $S_{\mathrm{wall}}$ 和随后的蓝队难度 $S_{\mathrm{bomb}}$。单纯增加搜索、延长探索、提高移动速度或覆盖更多轨迹，都不保证提高式\eqref{eq:red-total}。优化必须围绕“预期终端总分的边际增益”进行。

\section{来源等级与建模边界}

本文采用四类标签：
\begin{itemize}
  \item \textbf{源事实（S）}：可由固定提交中的文档或执行代码直接确认；
  \item \textbf{形式化（F）}：对源事实的等价数学表达，不增加行为假设；
  \item \textbf{近似（A）}：为求解而引入的有限栅格、独立性、线性化或候选集假设；
  \item \textbf{待标定（U）}：地图分布、参考蓝队族、运行环境和噪声联合分布中源码未完全指定的量。
\end{itemize}

真实墙体、真实位姿和评分产物只能作为模拟器隐状态与离线评分真值。比赛模式禁止控制器在函数体内调用真值 API，提交检查器也拒绝这些调用。\src{S3,S4} 因而，数学模型可以用真值描述系统转移和评分，但可部署策略必须只对本机公开观测、公开场地元数据、自己的动作历史和正式收到的消息可测。

\section{参数、坐标和索引}

设地图尺寸为 $W\times H$ 世界像素，世界原点在地图中心，$x$ 向右、$y$ 向上。提交栅格的原点在左上，转换为\src{S2}
\begin{equation}
c=\operatorname{round}(x+W/2),\qquad
r=\operatorname{round}(H/2-y).
\label{eq:world-grid}
\end{equation}
设红队无人机集合 $\mathcal I_R=\{1,\ldots,n_R\}$，蓝队集合 $\mathcal I_B=\{1,\ldots,n_B\}$；每架红机初始炸弹库存为 $k_i$，总参考数量 $K=\sum_i k_i$。这些是评估配置参数，不应无条件固定成某个数值。设仿真步上限为 $T$，墙钟上限为 $T^{\mathrm{wall}}$。默认红队配置示例有 $n_R=5,k_i=2$，但模型保留一般形式。\src{S5}

真墙栅格记为 $T_g\in\{0,1\}^{H\times W}$，其中 1 表示真实墙体实体栅格化后的占用。处置中心、返回区和干扰区不属于 $T_g$。边界墙掩码为 $B_g\in\{0,1\}^{H\times W}$。

\section{精确随机动态系统}

\subsection{全局隐状态}

在步 $t$，一个足够的 Markov 状态可写为
\begin{align}
X_t=(&M,\{p_{i,t},\theta_{i,t},v_{i,t},\omega_{i,t},h_{i,t},q_{i,t},g_{i,t},e_{i,t},b^{p}_{i,t-1}\}_{i\in\mathcal I_R\cup\mathcal I_B},\nonumber\\
&\mathcal B_t,\mathcal Z_t,\eta_t,\mathcal C_t,t,\tau^{\mathrm{wall}}_t).
\label{eq:state}
\end{align}
这里 $M$ 包含地图几何和公开场地参数；$p,\theta,v,\omega$ 是真实运动状态；$h$ 是健康度；$q$ 是红队剩余库存；$g$ 是抓取状态；$e$ 是设备可用状态；$b^p_{i,t-1}$ 是上一拍放置命令；$\mathcal B_t$ 是炸弹位置、速度、抓持和处置状态；$\mathcal Z_t$ 是区域状态；$\eta_t$ 是有色传感器噪声内部状态；$\mathcal C_t$ 是通信器缓存；$\tau^{\mathrm{wall}}_t$ 是墙钟时间。真值属于模拟器，不属于控制器观测。

源码对应的精确一步转移统一记为
\begin{equation}
X_{t+1}=\Phi(X_t,U_t,\xi_{t+1};\vartheta),
\label{eq:transition}
\end{equation}
其中 $\Phi$ 包含 Pymunk 刚体接触、阻尼、十个微步、设备更新和交互；$\xi$ 汇总随机噪声，$\vartheta$ 是固定引擎参数。采用 $\Phi$ 而不强行写成简单欧拉方程，是因为碰撞、抓取关节、炸弹载荷和精灵凸包几何不能由一个无接触线性模型精确替代。

可解释的自由运动部分为：无人机质量 50，本体系前向/横向命令经裁剪后，二维合力若范数超过 1 会再归一化；局部力标度为 $100\times3=300$，旋转命令直接设定角速度标度 $0.3\times0.6=0.18$。Playground 每次宏步执行 10 个 $0.1$ 的物理微步，空间阻尼为 0.95。\src{S11,S12} 这些常数解释动作量级，但不应被误写成“每步必然移动固定距离”。

\subsection{健康和区域}

健康初值为 50。碰墙与碰机共用一个基于 \texttt{time.perf\_counter()} 的计时器：距离上次扣血超过 1 秒墙钟时间后，相应碰撞才扣 1 点；健康不大于 0 时释放抓取并移除无人机。\src{S13,S14} 因而计算开销不仅消耗墙钟预算，还能改变碰撞扣血时序，这是源码事实，不是理想化仿真时间模型。

NO\_GPS 区会禁用 GPS 和罗盘，NO\_COM 区禁用通信器，KILL 区禁用设备并改变显示状态。\src{S15} 文档声称返回区逐渐回血，但对固定提交的运行时代码检索没有发现回血更新；精确模型不加入回血转移，只把此项列为文档与代码冲突。

\section{观测模型}

每架无人机 $i$ 的公开观测为
\begin{equation}
O_{i,t}=\mathcal O_i(X_t,\nu_{i,t}),
\end{equation}
包括可用时的 GPS、罗盘、里程计、Lidar、语义传感器、健康、库存、抓持、时间、场地尺寸以及上一步收到的正式消息。

\subsection{GPS 与罗盘的 AR(1) 噪声}

GPS 两轴测量为
\begin{align}
z^p_{i,t}&=p_{i,t}+a^p_{i,t},\\
a^p_{i,t}&=0.98a^p_{i,t-1}+\epsilon^p_{i,t},\qquad
\epsilon^p_{i,t}\sim\mathcal N\!\left(0,5^2(1-0.98^2)I_2\right).
\end{align}
罗盘为
\begin{align}
z^\theta_{i,t}&=\wrap(\theta_{i,t}+a^\theta_{i,t}),\\
a^\theta_{i,t}&=0.98a^\theta_{i,t-1}+\epsilon^\theta_{i,t},\qquad
\epsilon^\theta_{i,t}\sim\mathcal N\!\left(0,(4^\circ)^2(1-0.98^2)\right).
\end{align}
因此 5 px 和 $4^\circ$ 是稳态标准差，不是每步白噪声标准差；AR 初值为零。设备禁用时返回 \texttt{None}。\src{S8,S9}

\subsection{里程计}

令 $\Delta p_{i,t}=p_{i,t}-p_{i,t-1}$，则原始里程计为
\begin{equation}
d_t=\|\Delta p_{i,t}\|,quad
\alpha_t=\wrap(\operatorname{atan2}(\Delta y,\Delta x)-\theta_{i,t-1}),quad
\Delta\theta_t=\wrap(\theta_{i,t}-\theta_{i,t-1}),
\end{equation}
三个输出分别加独立高斯噪声，标准差为 $0.2$ px、$8^\circ$、$1^\circ$，角度再归一化。\src{S8}

公开辅助函数 \texttt{measured\_velocity()} 实际计算
\begin{equation}
\widehat v_t=\tilde d_t[\cos(z^\theta_t+\tilde\alpha_t),\ \sin(z^\theta_t+\tilde\alpha_t)]^\top.
\end{equation}
但 $\alpha_t$ 是相对于上一拍朝向定义的，辅助函数使用当前罗盘角且未减去 $\Delta\theta_t$；它是公开 API 的测量量，不应自动等同于真实速度或严格的位移向量。\src{S16}

\subsection{射线传感器}

Lidar 有 181 条、$360^\circ$、量程 300 px 的射线，距离加入 $\mathcal N(0,2.5^2)$ 噪声。语义传感器有 35 条、$360^\circ$、量程 200 px 的射线，返回命中实体类别、距离、相对角和抓持状态。\src{S2,S17,S18}

语义距离实现中调用的是 \texttt{np.random.normal(self.\_std\_dev\_noise)}。按 NumPy 的位置参数语义，这等价于 $\mathcal N(2.5,1^2)$，而不是看起来可能想表达的 $\mathcal N(0,2.5^2)$。\src{S18} 本模型在“源码精确模式”使用前者；若以后修正框架，必须把它作为模型版本变化，而不是在控制器里暗自假定已修正。

\section{动作、放置和通信约束}

每架无人机动作
\begin{equation}
u_{i,t}=(f_{i,t},\ell_{i,t},r_{i,t},g_{i,t},b_{i,t}),
\end{equation}
其中 $f,\ell,r\in[-1,1]$，$g\in\{0,1\}$；红队还有 $b\in\{0,1\}$ 的放置命令。\src{S2}

一次成功放置必须同时满足\src{S19}
\begin{equation}
b_{i,t}=1,\quad b_{i,t-1}=0,\quad q_{i,t}>0,\quad e^{\mathrm{deploy}}_{i,t}=1,
\end{equation}
并且没有抓住炸弹、抓取器不处于可抓状态、处置中心存在，且真实当前位置到每枚已有炸弹的距离均不少于 24 px。炸弹在\textbf{真实物理位置} $p_{i,t}$ 生成，库存只在成功时减一；失败仍会更新上一拍命令，所以再次尝试前必须先发 0。控制器能决定的是移动与触发时机，不能直接指定绝对落点。

通信半径为 250 px，收发双方都必须可用，实际距离判定是严格小于各自范围。通信不做墙体遮挡；消息可以是 Python 对象。每步先从上一拍接收缓存读取消息并产生本拍广播与控制，若没有提交，再经过物理步传递消息，供下一次控制读取。\src{S20,S21} 于是通信图为
\begin{equation}
\mathcal G_t=(\mathcal I,E_t),\qquad
(i,j)\in E_t\Longleftrightarrow \|p_{i,t}-p_{j,t}\|<250
\text{ 且双方通信器可用}.
\end{equation}

\section{信息结构与可部署策略}

无人机 $i$ 的本地历史定义为
\begin{equation}
H_{i,t}=(O_{i,0:t},U_{i,0:t-1},M^{\mathrm{recv}}_{i,0:t}),
\end{equation}
可部署策略必须满足
\begin{equation}
U_{i,t}\sim\pi_i(\cdot\mid H_{i,t}),\qquad
M^{\mathrm{send}}_{i,t}=\mu_i(H_{i,t}),
\label{eq:local-policy}
\end{equation}
且提交决定与提交栅格也只能由某台无人机的本地历史生成。任何跨实例共享全局状态、真实地图、真实位姿或离线评分产物都不在可行策略集合内。

集中式联合信念
\begin{equation}
b_t(x)=\Prob(X_t=x\mid H_{1,t},\ldots,H_{n,t})
\end{equation}
是分析上界，不是可直接部署的输入。由于多机消息可能转发同一观测、同一帧射线共享定位偏差，信念更新必须保存来源和相关性；把所有墙点当独立同分布观测会重复计数并低估位姿不确定性。

\section{红队终止事件}

红队停止时间是
\begin{equation}
\tau_R=\min\{\tau_{\mathrm{sub}},T,T^{\mathrm{wall}},\tau_{\mathrm{crash}}\}.
\end{equation}
任意一台无人机调用 \texttt{submit\_exploration\_map} 后，回合在下一次物理更新前结束；同一控制拍计算出的运动、放置和消息都不会执行。\src{S2,S21} 若多台同时提交，Launcher 按无人机遍历顺序取第一个非空栅格。提交不是“附带动作”，而是团队不可逆停止动作。

\section{地图精度分的精确公式}

令 $P\in\{0,1\}^{H\times W}$ 为提交栅格。对真实墙栅格 $T_g$ 计算每个墙像素到自由空间（包括栅格外部）的欧氏深度 $\delta_T(u)$，定义
\begin{equation}
A=\{u:T_g(u)=1,\ \delta_T(u)\le4\},\qquad
I=\{u:T_g(u)=1,\ \delta_T(u)>4\}.
\end{equation}
默认只给表面 $A$ 计分，深层内部权重为 0。命中信用为
\begin{equation}
C(P,T_g)=\frac{|\{u\in A:P(u)=1\}|}{|A|},
\label{eq:map-credit}
\end{equation}
若相应区域为空则按实现对可用区域重归一；默认场景可用式\eqref{eq:map-credit}。

令 $d_T(u)$ 为预测像素 $u$ 到最近真实墙像素的精确欧氏距离。每个预测像素的收费为
\begin{equation}
c(d)=
\begin{cases}
0,&d\le10,\\
\min\left(1,0.1\left\lceil\dfrac{d-10}{2}\right\rceil\right),&d>10.
\end{cases}
\end{equation}
总扣分
\begin{equation}
Q(P,T_g)=\frac{0.2}{|A|}\sum_{u:P(u)=1}c(d_T(u)),
\end{equation}
最终
\begin{equation}
S_{\mathrm{map}}=100\clip(C-Q,0,1).
\label{eq:map-score}
\end{equation}
关键含义是：10 px 免罚区只免除\textbf{误画扣分}，并不会让距离真墙 10 px 内的预测像素获得命中信用。命中仍要求在同一真墙像素上 $P(u)=1$。地图内部的时间诊断权重为 0，时间只通过顶层 $S_{\mathrm{time}}$ 计入。\src{S1,S6}

\section{炸墙覆盖分的精确公式}

默认爆炸半径像素数
\begin{equation}
R=\operatorname{round}(0.12\min(H,W)).
\end{equation}
内部墙像素权重为 1，边界框墙像素权重为 0.5：
\begin{equation}
w(u)=T_g(u)\,[1-0.5B_g(u)].
\end{equation}
对实际成功放置的中心集合 $\mathcal P$，圆盘并集
\begin{equation}
\mathcal D(\mathcal P)=\bigcup_{p\in\mathcal P}\{u:\|u-p\|_2\le R\}
\end{equation}
产生实际加权覆盖
\begin{equation}
D(\mathcal P,T_g)=\sum_u w(u)\one\{u\in\mathcal D(\mathcal P)\}.
\end{equation}
重叠区域不重复计数。

分母不是所有墙的总权重，而是一个确定性贪心参考 $G_K(T_g)$。从所有自由像素候选中心开始，连续 $K$ 次选择新增未覆盖加权墙体最多的位置，\texttt{np.argmax} 以行优先第一个候选打破平局，没有正增益时提前停止。于是
\begin{equation}
S_{\mathrm{wall}}=100\clip\left(\frac{D(\mathcal P,T_g)}{G_K(T_g)},0,1\right).
\label{eq:wall-score}
\end{equation}
\src{S1}

需要特别说明：$G_K$ 是栅格自由点上的贪心松弛参考，不是带运动学、机体碰撞、有限时间、24 px 炸弹间距和信息约束的物理全局最优。源码注释把它称为可达参考，但其候选筛选没有编码这些全部约束。实际圆盘用距离变换判定 $d\le R$；贪心卷积核使用 $x^2+y^2\le R^2+0.5$，边缘上有极小的离散化差异。本文保留源码口径，不把 $G_K$ 冒称严格物理最优。

\section{时间分}

令默认最佳比例 $\beta=0.5$，红队实际结束步为 $\tau_R$。无崩溃时
\begin{equation}
S_{\mathrm{time}}=100\clip\left(\frac{T-\tau_R}{T-\beta T},0,1\right).
\label{eq:time-score}
\end{equation}
即不超过一半时限均为 100，之后线性降至 0；崩溃直接为 0。默认使用仿真步数而非墙钟时间。\src{S1,S2} 墙钟仍是硬终止约束，所以“时间分满分”不代表计算成本无关：过慢的控制器可能先触发墙钟上限，也会影响基于墙钟的碰撞扣血。

\section{蓝队过程与炸弹难度分}

红队提交的实际炸弹布局 $\mathcal P$ 进入参考蓝队救援回合。设布局炸弹数 $N_P$，成功处置数 $N_D$，蓝队初始无人机数 $n_B$，在返回区内（若无返回区则全场）的健康总和为 $H_R$，蓝队结束步为 $\tau_B$。蓝队分数为\src{S2,S7}
\begin{align}
S_B&=0.6D_B+0.2H_B+0.2T_B,\\
D_B&=\begin{cases}100N_D/N_P,&N_P>0,\\100,&N_P=0,\end{cases}\\
H_B&=100\frac{H_R}{50n_B},\\
T_B&=100\clip\left(\frac{T_B^{\max}-\tau_B}{T_B^{\max}-0.5T_B^{\max}},0,1\right),
\end{align}
崩溃时蓝队时间分为 0。全部炸弹处置完成会提前终止；未完成通常运行到限制。健康分的分母始终包含全部初始蓝机，所以没有返回的蓝机贡献为零，而不是只对返回者求平均。

第 $j$ 个有效参考蓝队运行的布局难度为
\begin{equation}
d_j=\clip(2(100-S_{B,j}),0,100).
\end{equation}
若 $J_{\mathrm{ok}}$ 个有效参考运行满足最低成功数量，默认聚合为均值
\begin{equation}
\bar d=\frac1{J_{\mathrm{ok}}}\sum_{j=1}^{J_{\mathrm{ok}}}d_j;
\end{equation}
可配置为中位数。失败或崩溃的参考运行被排除；若有效数量不足，炸弹分为 0。红队成功放置数为 $N_P$、期望数为 $K$ 时
\begin{equation}
S_{\mathrm{bomb}}=\min(1,N_P/K)\,\bar d.
\label{eq:bomb-score}
\end{equation}
\src{S1} 这说明红蓝关系不是零和：蓝队分仅通过饱和变换影响红队 20\% 的一项，而地图、炸墙和时间仍独立存在。

\section{完整优化问题}

设评测地图与配置 $m\sim\mu_M$，模拟随机性 $\xi\sim\mu_\xi$，参考蓝队策略索引 $j\sim\mu_B$。红队策略族 $\Pi$ 必须满足式\eqref{eq:local-policy}的信息约束、提交限制和打包规则。完整目标为
\begin{equation}
\max_{\pi\in\Pi}\ J(\pi)
=\E_{m,\xi,j}\!\left[
0.5S_{\mathrm{map}}+0.2S_{\mathrm{bomb}}+0.2S_{\mathrm{wall}}+0.1S_{\mathrm{time}}
\right].
\label{eq:full-objective}
\end{equation}
如果正式赛制给每个配置权重 $\lambda_c$，则外层期望应使用
\begin{equation}
J(\pi)=\frac{\sum_c\lambda_c\,\E[S_R\mid c]}{\sum_c\lambda_c}.
\end{equation}
“先回合平均、再种子平均、再地图等权平均”是对某个特定 $\mu_M,\mu_\xi$ 的样本估计器，不是源码内在规则。原始五图或新增地图如何加权，必须由实际比赛评估计划决定。

\subsection{随机种子的正确角色}

对候选策略 $a,b$，合理的快速筛选是使用共同随机数作配对差分
\begin{equation}
\widehat\Delta=\frac1N\sum_{r=1}^{N}[S_R(a;\omega_r)-S_R(b;\omega_r)],
\end{equation}
再报告均值、最低分和差分不确定度。种子不进入策略逻辑，也不能按种子硬编码目标。由于固定提交没有统一播种所有随机源，相同数值种子仍不保证 $\omega_r$ 完全相同；真正的可重复评测需要框架统一管理所有随机生成器，或者接受统计重复而非逐位重复。

\section{探索、放置、继续和提交的动态决策}

在集中式信念上界中，剩余 $r$ 步、信念 $b$、库存 $q$ 的 Bellman 方程为
\begin{equation}
V_r(b,q)=\max\left\{
Q_{\mathrm{sub}}(b,q),
\sup_{u\in\mathcal U_{\mathrm{move}}}Q_{\mathrm{move}}(b,q,u),
\sup_{u\in\mathcal U_{\mathrm{place}}}Q_{\mathrm{place}}(b,q,u)
\right\},
\label{eq:bellman}
\end{equation}
其中
\begin{align}
Q_{\mathrm{sub}}(b,q)&=\E[S_R\mid b,\text{立即提交}],\\
Q_a(b,q)&=-c_{\mathrm{wallclock}}(a)-c_{\mathrm{risk}}(a)
+\E[V_{r-1}(b',q')\mid b,a].
\end{align}
放置的边际价值必须同时包含
\begin{equation}
\Delta_{\mathrm{place}}
=0.2\Delta S_{\mathrm{wall}}+0.2\Delta S_{\mathrm{bomb}}
-\Delta C_{\mathrm{inventory}}-\Delta C_{\mathrm{travel}}-\Delta C_{\mathrm{risk}}-\Delta C_{\mathrm{compute}}.
\end{equation}
探索的边际价值则主要通过降低地图和可达落点信念的不确定度，改变未来 $S_{\mathrm{map}},S_{\mathrm{wall}},S_{\mathrm{bomb}}$。提交在
\begin{equation}
Q_{\mathrm{sub}}\ge\max(Q_{\mathrm{move}},Q_{\mathrm{place}})
\end{equation}
时才合理。这一准则自然随已知墙体、位姿不确定度、剩余库存、剩余时间、可通信性和蓝队布局难度变化，不需要固定时间阈值加固定库存保留。

式\eqref{eq:bellman}是分析上界。精确分散求解应使用公共信息状态 $\Pi_t$（对各机私人历史的条件分布）和对本地历史的处方 $\gamma_{i,t}:H_{i,t}\mapsto U_{i,t}$；实际工程可用局部信念、带来源的地图摘要和邻接通信近似，但必须承认由此产生的最优性缺口。

\section{已知信念下的有限候选近似}

当需要把布局规划转成可计算问题时，可从当前后验生成有限候选落点 $\mathcal L$ 和可行移动图 $E$。令 $x_{i\ell}\in\{0,1\}$ 表示无人机 $i$ 是否在候选 $\ell$ 放置，$y_u\in\{0,1\}$ 表示墙像素 $u$ 是否被至少一个选中圆盘覆盖。一个基本近似为
\begin{align}
\max_x\quad &0.2\sum_u \widehat w_u y_u
+0.2\widehat D_{\mathrm{blue}}(x)
-\lambda_t C_{\mathrm{route}}(x)-\lambda_r C_{\mathrm{risk}}(x),\\
\text{s.t.}\quad
&y_u\le\sum_{i,\ell:u\in\mathcal D(\ell)}x_{i\ell},\\
&\sum_\ell x_{i\ell}\le q_i,\\
&x_{i\ell}+x_{j\ell'}\le1\quad\text{若 }\|\ell-\ell'\|<24,\\
&\text{流守恒、到达时刻、视线安全裕度与剩余时域约束。}
\end{align}
这里 $\widehat w_u=\E[w(u)\mid b_t]$，$\widehat D_{\mathrm{blue}}$ 必须由参考蓝队滚动仿真或经过校准的代理模型给出；只用到处置中心或出生点的距离替代它，不是源代码等价公式。这个模型是候选离散化近似，不是完整游戏的精确解。

\section{分数门槛的代数约束}

令四个归一化分项为 $m,d,w,t\in[0,1]$。总分不低于 95 等价于损失预算
\begin{equation}
50(1-m)+20(1-d)+20(1-w)+10(1-t)\le5.
\label{eq:loss-budget}
\end{equation}
由此得到单项必要条件
\begin{equation}
m\ge0.90,\qquad d\ge0.75,
\qquad w\ge0.75,qquad t\ge0.50,
\end{equation}
但这些条件联合起来并不充分。例如每项恰好达到该下界会损失 20 分。正确用法是按式\eqref{eq:loss-budget}分配总损失，而不是分别追求看似不错的阈值。

若缺少一枚炸弹且其余条件不变，式\eqref{eq:bomb-score}的完成率最多下降 $1/K$，仅炸弹难度项最多损失 $20/K$ 分；同时它还可能降低 $S_{\mathrm{wall}}$，所以真实损失通常更大。时间在前半时域为平台：只要墙钟安全裕度足够，在 $T/2$ 前用额外步骤改善地图、炸墙或蓝队难度，没有顶层时间分代价。

\section{模块有效性的数学审查框架}

以后每个模块不能只看“被调用了多少次”，应记录五元组
\begin{equation}
\mathcal A_k=(N_k,C_k,\Delta U_k,\Delta Z_k,\Delta S_k),
\end{equation}
其中 $N_k$ 是完整调用次数，$C_k$ 是累计与尾部耗时，$\Delta U_k$ 是对最终动作的实际扰动，$\Delta Z_k$ 是对中间信念、地图、路线或消息的扰动，$\Delta S_k$ 是在配对样本上的因果分数变化。按证据分为：未调用；调用但输出被覆盖或无决策影响；有决策影响但未证实改善；配对评测证实有效。Top-N profile 中未出现不能当作零调用，单图零调用也不能推断永久无效。

缓存模块还需报告请求数、命中数、失效原因、节省的底层计算次数和额外维护时间；搜索模块需报告候选生成数、完整评估数、接受数和被各约束拒绝数。一个模块的保留条件应是
\begin{equation}
\E[\Delta S_k]-\lambda_C\E[C_k]-\lambda_R\operatorname{Risk}_k>0,
\end{equation}
而不是代码行数多、调用次数多或局部微基准更快。

\section{资料冲突和不可静默填补的未知量}

\begin{longtable}{p{0.25\textwidth}p{0.67\textwidth}}
\toprule
事项 & 源审计结论\\
\midrule
返回区回血 & 文档称逐渐恢复，但固定提交运行时代码没有找到健康增加路径。精确模型不加入回血；需要上游确认或修改框架后重新版本化。\\
语义距离噪声 & 实现的单位置参数调用对应 $\mathcal N(2.5,1)$；若意图是零均值标准差 2.5，文档/实现不一致。\\
辅助速度语义 & 文档称像素/秒，实际由单步里程和当前罗盘构造，且参考角定义存在一拍差异；应作为测量特征，不当作真值。\\
炸墙“理论最优” & 实际是确定性贪心栅格参考，未编码全部物理可达约束；它不是可证明的连续物理全局最优。\\
随机种子 & 原资料没有三固定种子规则；多个随机源也没有统一由 Playground 种子控制。\\
地图分布与正式权重 & 固定提交给出评估计划机制和示例配置，但未来新增地图及其正式权重属于外部赛制输入。\\
参考蓝队族 & 炸弹难度依赖实际被采用的参考蓝队策略与有效运行集合；不能只从红队布局解析式得到。\\
\bottomrule
\end{longtable}

\section{验证清单}

数学模型可以按以下层次验证：
\begin{enumerate}
  \item \textbf{公式单元验证}：用小型二值墙图手工核对命中、免罚、阶梯扣分、圆盘并集、边界权重、贪心分母和时间平台；
  \item \textbf{源码对照验证}：对保存的评分输入调用原始 \texttt{PlaceScoreManager}，逐项比对公式输出，而不是只比最终总分；
  \item \textbf{信息约束验证}：检查策略每个输入的来源，证明没有真值、地图类、评估产物或跨实例状态；
  \item \textbf{动力学验证}：用公开传感器和命令记录拟合局部转移，只把它当工作点模型，不替换含碰撞的 $\Phi$；
  \item \textbf{策略因果验证}：在共同样本上做模块消融，报告四个分项、最低回合、运行时间和置信区间；
  \item \textbf{阶段门槛验证}：当前人工规则“Map01--Map05 每图均分均达到 95 后再扩图”应逐图报告，不能由五图总平均替代。这个门槛是用户的研发流程，不是 Gitee 游戏规则。
\end{enumerate}

\section{对后续算法设计的直接约束}

从模型可直接得出以下设计原则。地图估计应保留位姿与墙体后验，并以式\eqref{eq:map-score}的期望增益而不是占用栅格数量为目标；放置规划必须同时评估式\eqref{eq:wall-score}和式\eqref{eq:bomb-score}，不能只最大化几何覆盖；提交必须视作停止问题，且提交拍不能再期待动作生效；通信摘要必须带来源与时间戳，避免重复融合；运行时间优化要计算到完整终端分数的机会成本，而不是只追求每步更快。

因此，下一轮算法工作应从三个可测对象开始：当前信念下的地图分期望损失、每个候选落点对墙分和蓝队难度的联合边际贡献、继续一拍与立即提交的价值差。任何新模块若不能改变这三个对象之一，或者改变了动作却没有配对分数证据，就不应继续消耗正式仿真预算。

\section{源文件索引}

下列编号均指固定提交 \texttt{53e129bd7d2efc82c29deac2ecd8c57a536e74e7}；文件哈希另见同目录 \path{source_manifest.json}。

\begin{longtable}{p{0.08\textwidth}p{0.84\textwidth}}
\toprule
编号 & 固定提交中的来源\\
\midrule
S1 & \path{src/swarm_rescue/simulation/reporting/place_score_manager.py}：红队权重、地图/炸弹/炸墙/时间分、贪心参考、合并总分。\\
S2 & \path{doc/participants/02-队伍开发指南.md}：公开 API、动作、提交格式、评分说明、蓝队规则。\\
S3 & \path{doc/participants/01-安装与配置指南.md}：比赛模式真值访问限制。\\
S4 & \path{doc/participants/03-提交检查与打包.md}：真值 API 与禁止目录的静态检查。\\
S5 & \path{config/competition_place_eval_plan.yml} 与 \path{config/competition_rescue_eval_plan.yml}：示例评估计划、机数、回合数和库存。\\
S6 & \path{src/swarm_rescue/tools/map_exploration_scoring.py}：墙面/内部拆分、像素命中、误画扣分和地图分。\\
S7 & \path{src/swarm_rescue/simulation/reporting/score_manager.py} 与 \path{simulation/gui_map/map_abstract.py}：蓝队三项分数与返航健康。\\
S8 & \path{src/swarm_rescue/simulation/drone/drone_sensors.py}：GPS、罗盘和里程计观测。\\
S9 & \path{src/swarm_rescue/simulation/utils/utils_noise.py}：高斯与 AR(1) 噪声更新。\\
S10 & \path{src/swarm_rescue/simulation/gui_map/playground.py}：独立随机生成器、物理步和消息传输。\\
S11 & \path{src/swarm_rescue/simulation/drone/drone_base.py}：质量、命令与力/角速度标度。\\
S12 & \path{src/swarm_rescue/simulation/utils/constants.py}、\path{definitions.py}：速度比例、力、角速度和传感器规格。\\
S13 & \path{src/swarm_rescue/simulation/drone/drone_abstract.py}：健康、碰撞、公开测量、真值保护。\\
S14 & \path{src/swarm_rescue/simulation/utils/timer.py}：基于 \texttt{perf\_counter} 的墙钟计时器。\\
S15 & \path{src/swarm_rescue/simulation/elements/sensor_disablers.py}：禁用区域行为。\\
S16 & \path{src/swarm_rescue/simulation/drone/drone_abstract.py}：\texttt{measured\_velocity} 与角速度辅助函数。\\
S17 & \path{src/swarm_rescue/simulation/ray_sensors/drone_lidar.py}：Lidar 量程、射线与噪声。\\
S18 & \path{src/swarm_rescue/simulation/ray_sensors/drone_semantic_sensor.py}：语义观测和距离噪声实现。\\
S19 & \path{src/swarm_rescue/simulation/drone/bomb_deployer.py}：库存、上升沿、24 px 间距和真实物理位置放置。\\
S20 & \path{src/swarm_rescue/simulation/drone/communicator.py}：通信范围、收发条件和缓存。\\
S21 & \path{src/swarm_rescue/simulation/gui_map/gui_sr.py}：控制、提交、物理、消息和终止的执行顺序。\\
\bottomrule
\end{longtable}

\section*{模型版本声明}

本版是“固定提交源码精确层 + 明示的优化近似层”。它完整定义了游戏状态、公开观测、动作、通信、终止、四个红队分项、蓝队终端过程、随机评测和可部署策略约束。它没有假装已知未来地图分布或参考蓝队的闭式行为；这些必须通过正式配置和实际蓝队运行补齐。若上游提交、评分配置或参考蓝队改变，应更新来源哈希并生成新的模型版本，不能把不同版本的实测分数混合解释。

\end{document}
